\section{The three-example CoHA $\to$ vertex-algebra $\to$ gauge chain, first-principles}
\label{w17:sec:coha-c3-to-w-infty-agent13}

\providecommand{\ClaimStatusTheorem}{\textsc{[theorem, cited]}}
\providecommand{\ClaimStatusConjectured}{\textsc{[conjectural]}}
\providecommand{\ClaimStatusDefinition}{\textsc{[definition]}}
\providecommand{\ClaimStatusDerivation}{\textsc{[first-principles derivation]}}
\providecommand{\ClaimStatusOpen}{\textsc{[open]}}
\providecommand{\ClaimStatusProvedElsewhere}{\textsc{[proved elsewhere]}}
\providecommand{\CoHA}{\mathrm{CoHA}}
\providecommand{\Sh}{\mathrm{Sh}}
\providecommand{\Wilpha}{\mathcal{W}_{1+\infty}}
\providecommand{\gghone}{\widehat{\mathfrak{gl}}_1}
\providecommand{\hCS}{\mathrm{hCS}}
\providecommand{\Jac}{J}
\providecommand{\Fact}{\mathrm{Fact}}
\providecommand{\Crit}{\mathrm{Crit}}
\providecommand{\Hilb}{\mathrm{Hilb}}
\providecommand{\Obs}{\mathrm{Obs}}
\providecommand{\eps}{\epsilon}
\providecommand{\Zinst}{Z^{\mathrm{inst}}_{\Omega}}
\providecommand{\PhiFA}{\Phi^{\mathrm{FA}}}
\providecommand{\SpCh}{\mathrm{Sp}^{\mathrm{ch}}}
\providecommand{\Mred}{\mathcal{M}^{\mathrm{red}}_{DT}}
\providecommand{\Auts}{\mathrm{Aut}_s}
\providecommand{\gBPS}{\mathfrak{g}_{\mathrm{BPS}}}
\providecommand{\gDelta}{\mathfrak{g}_{\Delta_5}}
\providecommand{\bY}{\mathbb{Y}}
\providecommand{\Qcon}{Q_{\mathrm{con}}}

\emph{Scope.} Three worked examples of the CY$_3$--CoHA $\to$ vertex-algebra $\to$
gauge-theory chain, at successively higher difficulty: $\mathbb{C}^3$ (flat, toric,
local, unconditional), the resolved conifold (non-flat, toric, local, partial), and
$K3 \times E$ (compact, non-toric, conjectural; graded-dimension matching only).
Each identification is pinned to a primary citation, an explicit chain-level
computation, and a gauge-theory incarnation where available. The three identifications
that must be kept separate --- $\CoHA(\C^3) \neq Y(\gghone)$, $Y(\gghone) \neq \Wilpha$,
and $\CoHA(\C^3) \neq \Wilpha$ --- are tagged with scope at every use. The treatise
cross-consistency target is \texttt{\textasciitilde /calabi-yau-quantum-groups/
notes/CoHA\_to\_W\_infty\_treatise.tex}: every theorem, remark, and open-status
claim here coincides with the treatise, and any disagreement is flagged explicitly.

%--------------------------------------------------------------------------------
\subsection{Example 1: $\mathbb{C}^3$ and the Jordan triple loop}
\label{w17:subsec:c3-full-chain}

\subsubsection{The quiver, its potential, and its Jacobi algebra}
\label{w17:subsubsec:jordan-triple-jacobi}

\begin{definition}[Jordan triple loop quiver $Q$]
\label{w17:def:jordan-triple-quiver}
\ClaimStatusDefinition\ The Jordan triple loop quiver has one vertex $*$ and three
oriented loops $X, Y, Z: * \to *$. The path algebra is the free associative
$k$-algebra $kQ = k\langle X, Y, Z\rangle$.
\end{definition}

\begin{definition}[Cubic cyclic potential]
\label{w17:def:cubic-potential}
\ClaimStatusDefinition\ The potential is the cyclic-invariant word
\[
 W \;:=\; \mathrm{tr}\bigl(X[Y,Z]\bigr) \;=\; \mathrm{tr}(XYZ - XZY)
 \;\in\; kQ_{\mathrm{cyc}} := kQ/[kQ, kQ].
\]
Cyclic invariance is automatic: $XYZ, YZX, ZXY$ represent the same class in
$kQ_{\mathrm{cyc}}$, and the sign flip under transposition gives $-XZY = -XZY$
stable under cyclic rotation of its three letters.
\end{definition}

\begin{proposition}[Jacobi algebra first-principles]
\label{w17:prop:jacobi-c3}
\ClaimStatusDerivation\
$\Jac(Q, W) := kQ/\langle \partial_a W : a \in \{X, Y, Z\}\rangle$ is the
polynomial algebra $k[X, Y, Z] = \mathcal{O}_{\C^3}$.
\end{proposition}

\begin{proof}
The cyclic derivative $\partial_a: kQ_{\mathrm{cyc}} \to kQ$ acts on a word
$m = a_1 a_2 \cdots a_n$ by
\[
 \partial_a m \;=\; \sum_{\substack{1 \leq i \leq n \\ a_i = a}}
 a_{i+1} a_{i+2} \cdots a_n a_1 a_2 \cdots a_{i-1},
\]
extended bilinearly (Ginzburg~2006 Def.~3.5.1). Applying to $W = XYZ - XZY$:

\emph{$\partial_X$:} letter $X$ appears at position~$1$ of $XYZ$ with cyclic
successor $YZ$; letter $X$ appears at position~$1$ of $XZY$ with cyclic successor
$ZY$. The sign from the $-$ coefficient of $XZY$ gives
\[
 \partial_X W \;=\; YZ - ZY \;=\; [Y, Z].
\]

\emph{$\partial_Y$:} cyclic rotation of $XYZ$ to $YZX$ places $Y$ at position~$1$
with successor $ZX$; rotation of $XZY$ to $YXZ$ places $Y$ at position~$1$ with
successor $XZ$. With the $-$ sign on $XZY$:
\[
 \partial_Y W \;=\; ZX - XZ \;=\; [Z, X].
\]

\emph{$\partial_Z$:} rotation $XYZ \to ZXY$ gives successor $XY$;
rotation $XZY \to ZYX$ gives successor $YX$. With the $-$ sign:
\[
 \partial_Z W \;=\; XY - YX \;=\; [X, Y].
\]

The three relations $\{[Y, Z] = 0, [Z, X] = 0, [X, Y] = 0\}$ generate the ideal
killing all three commutators, so
\[
 \Jac(Q, W) \;=\; k\langle X, Y, Z\rangle / \langle [X, Y], [Y, Z], [Z, X]\rangle
 \;=\; k[X, Y, Z] \;=\; \mathcal{O}_{\C^3}.
\]
\end{proof}

\begin{verification}[Three paths to Proposition~\ref{w17:prop:jacobi-c3}]
\label{w17:verif:jacobi-c3}\mbox{}
\begin{enumerate}
\item \emph{Direct cyclic-derivative expansion.} Above.
\item \emph{CY$_3$-algebra characterisation.} The Ginzburg dg-algebra
$\Gamma(Q, W)$ has $H^0(\Gamma) = \Jac(Q, W)$ and carries a canonical
3-Calabi--Yau structure (Ginzburg~2006 Thm.~3.6.4). Computing
$\mathrm{Ext}^\bullet_{\mathcal{O}_{\C^3}}(k_0, k_0)$ of the origin skyscraper
sheaf directly:
$\mathrm{Ext}^i = \wedge^i (T_{\C^3, 0})^\vee$, a Koszul complex with generator
in degree $3$ (the top exterior power of $T^\vee_{\C^3, 0} \cong k\{dX, dY, dZ\}$);
this matches the CY$_3$ prediction.
\item \emph{Moduli of representations.} $\mathrm{Rep}_n(\Jac(Q, W)) = \{(X, Y, Z) \in
\mathrm{End}(\C^n)^3 : [X, Y] = [Y, Z] = [Z, X] = 0\}$ is the commuting variety.
After GIT-quotient by $GL_n$ with an ample stability condition, this recovers
$\Hilb^n(\C^3)$ (the Hilbert scheme of $n$ points on $\C^3$), which parametrises
colength-$n$ ideals of $\mathcal{O}_{\C^3}$ (Nakajima~1999; MNOP~2006). The
bijection between commuting triples and colength-$n$ ideals is an independent
check that $\Jac(Q, W) = \mathcal{O}_{\C^3}$.
\end{enumerate}
\end{verification}

\begin{remark}[Sign convention: cyclic vs.\ non-cyclic derivative]
\label{w17:rem:sign-convention}
The relations above use the cyclic derivative on $kQ_{\mathrm{cyc}}$, not the
ordinary (non-cyclic) derivative on $kQ$. For a cyclically invariant element,
the two agree up to a factor equal to the length of the word, but the
non-cyclic variant double-counts and gives $YZ - YZ = 0$ (and similarly for
the other letters), missing the Jacobi relations entirely. The cyclic variant
is uniquely determined by requiring $\partial_a(b) = \delta_{ab}$ on generators
and compatibility with cyclic rotation (Ginzburg~2006 Def.~3.5.1); no convention
ambiguity remains.
\end{remark}

\subsubsection{Moduli stack, torus action, CY$_3$ slice}
\label{w17:subsubsec:moduli-torus-cy3}

\begin{definition}[Representation moduli stack]
\label{w17:def:moduli-cn}
\ClaimStatusDefinition\ For dimension $n$,
\[
 \mathcal{M}_n \;:=\; \bigl[\{(X, Y, Z) \in \mathrm{End}(\C^n)^3 :
 [X, Y] = [Y, Z] = [Z, X] = 0\}/GL_n(\C)\bigr].
\]
\end{definition}

\begin{definition}[Torus action and CY$_3$ slice]
\label{w17:def:torus-cy3-slice}
\ClaimStatusDefinition\ $T := (\C^\times)^3$ acts on $(X, Y, Z)$ by
\[
 (t_1, t_2, t_3) \cdot (X, Y, Z) \;=\; (t_1 X, t_2 Y, t_3 Z).
\]
The CY$_3$ volume form $\Omega_{\C^3} = dX \wedge dY \wedge dZ$ has weight
$t_1 t_2 t_3$; it is $T$-invariant iff $t_1 t_2 t_3 = 1$, i.e.\ writing
$t_i = e^{\eps_i}$:
\[
 \boxed{\,\eps_1 + \eps_2 + \eps_3 \;=\; 0\,}
\]
(the CY$_3$ slice of the $T$-weight space). The equivariant parameter ring is
$\mathbb{F} := \C(\eps_1, \eps_2)$, with $\eps_3 = -\eps_1 - \eps_2$ on the slice.
\end{definition}

\begin{lemma}[Plane-partition parametrisation of $T$-fixed ideals]
\label{w17:lem:plane-partitions}
\ClaimStatusDerivation\ The $T$-fixed locus of $\Hilb^n(\C^3)$ is in bijection with
plane partitions (3d Young diagrams) of total volume $n$.
\end{lemma}

\begin{proof}
A $T$-fixed ideal $I \subset k[X, Y, Z]$ is $T$-invariant, hence generated by
$T$-eigenvectors. These are exactly the monomials $X^a Y^b Z^c$ of tri-weight
$(a\eps_1 + b\eps_2 + c\eps_3)$. Hence $I$ is a monomial ideal, determined by the
subset $S \subset \Z_{\geq 0}^3$ of monomials \emph{not} in $I$. Since $I$ is an
ideal, $S$ satisfies the downward-closure property: if $X^a Y^b Z^c \in S$
(i.e.\ $X^a Y^b Z^c \notin I$), then so are all $X^{a'} Y^{b'} Z^{c'}$ with
$0 \leq a' \leq a$, $0 \leq b' \leq b$, $0 \leq c' \leq c$.

This is the definition of a 3d Young diagram (plane partition). Colength $n$
means $|S| = n$. This is MNOP~2006 Prop.~5.1
(\texttt{arXiv:math/0312059}), with the combinatorics tracing back to
Cheah~1996.
\end{proof}

\begin{verification}[Three paths to Lemma~\ref{w17:lem:plane-partitions}]
\label{w17:verif:plane-partitions}\mbox{}
\begin{enumerate}
\item \emph{Direct combinatorial argument.} Above.
\item \emph{MacMahon generating function.} The generating function
$\sum_{n \geq 0} q^n |\Hilb^n(\C^3)^T| = M(q) := \prod_{k \geq 1}(1 - q^k)^{-k}$
(MacMahon 3d-partition count). Nekrasov~2003 \texttt{arXiv:hep-th/0206161}
identifies $M(q)$ with the $U(1)$ DT generating function, providing an
independent numerical check.
\item \emph{Equivariant character at a fixed point.} The $T$-character of the
tangent space $T_{[I]}\Hilb^n(\C^3)$ at a monomial-ideal fixed point is
computable via Okounkov~2003 \texttt{arXiv:math/0512330} Prop.~3 and matches
the plane-partition character formula by independent localisation argument.
\end{enumerate}
\end{verification}

\subsubsection{The CoHA of $\C^3$ and the Schiffmann--Vasserot shuffle}
\label{w17:subsubsec:coha-shuffle}

\begin{definition}[Cohomological Hall algebra with potential]
\label{w17:def:coha}
\ClaimStatusDefinition\ (Kontsevich--Soibelman~2008 \texttt{arXiv:0811.2435} \S6.)
\[
 \CoHA(\C^3) \;:=\; \bigoplus_{n \geq 0} H^*_T\bigl(\mathcal{M}_n,\, \phi_{W_n}\bigr)
\]
where $W_n = \mathrm{tr}_n(W): \mathrm{End}(\C^n)^3 \to \C$ is the trace of the
potential on the representation space, and $\phi_{W_n}$ is the vanishing-cycle
sheaf (Deligne SGA~7.2 Exp.~XIV). The multiplication is via extension correspondences:
for a short exact sequence $0 \to V' \to V \to V'' \to 0$ of commuting triples, one
pulls back a class on $\mathcal{M}_{n'} \times \mathcal{M}_{n''}$ and pushes forward to
$\mathcal{M}_{n}$ along the projection from the stack of flags.
\end{definition}

\begin{remark}[Vanishing cycle localisation to the commuting locus]
\label{w17:rem:vanishing-cycle-local}
The critical locus of $W_n$ on $\mathrm{End}(\C^n)^3$ is precisely the commuting
variety: $\Crit(W_n) = \{[X, Y] = [Y, Z] = [Z, X] = 0\}$ by the computation of
$\partial_{X_{ij}} W_n$ in Proposition~\ref{w17:prop:jacobi-c3} applied matrix-by-matrix.
Under mild regularity hypotheses (Schiffmann--Vasserot~2013
\texttt{arXiv:1202.2756} Thm.~1.1; Davison~2012 for general quiver-with-potential),
\[
 H^*_T(\mathrm{End}(\C^n)^3, \phi_{W_n})
 \;\simeq\; H^*_T(\Crit(W_n), \C)
\]
up to a cohomological shift. $T$-localisation then reduces to a sum over plane
partitions.
\end{remark}

\begin{theorem}[Schiffmann--Vasserot shuffle realisation]
\label{w17:thm:sv-shuffle}
\ClaimStatusProvedElsewhere\ (Schiffmann--Vasserot~2013 \texttt{arXiv:1202.2756}
Thm.~1.1.) Over the field $\mathbb{F}(\!(z)\!)$,
\[
 \CoHA(\C^3) \otimes_{\mathbb{F}} \mathbb{F}(\!(z)\!) \;\cong\; \Sh,
\]
where $\Sh := \bigoplus_{m \geq 0} \mathbb{F}[z_1, \ldots, z_m]^{S_m}$ is the
shuffle algebra with product
\begin{equation}
\label{w17:eq:shuffle-product}
 (F \star G)(x_1, \ldots, x_{m+n})
 \;:=\; \frac{1}{m!\, n!}
 \sum_{\sigma \in S_{m+n}} \sigma \cdot \Bigl[
  F(x_1, \ldots, x_m)\, G(x_{m+1}, \ldots, x_{m+n})
  \prod_{\substack{1 \leq i \leq m \\ m < j \leq m+n}} \omega(x_j, x_i)
 \Bigr]
\end{equation}
and kernel
\begin{equation}
\label{w17:eq:shuffle-kernel}
 \omega(z, w) \;:=\;
 \frac{(z - w - \eps_1)(z - w - \eps_2)(z - w - \eps_3)}{(z - w)^3}.
\end{equation}
\end{theorem}

\begin{proof}[Sketch, following Schiffmann--Vasserot~2013]
The proof routes via $T$-localisation: the integration map
$H^*_T(\mathcal{M}_n, \phi_{W_n}) \to H^*_T(\mathcal{M}_n^T, \phi_{W_n}^T)$ becomes
an isomorphism after inverting the equivariant parameters, with localisation giving
an explicit plane-partition contribution per fixed point. The Hall correspondence
product translates, after localisation, into the symmetric-group-symmetrised
combinatorial shuffle~\eqref{w17:eq:shuffle-product}. The kernel
$\omega(z, w)$ encodes the virtual-tangent $T$-character
\[
 \mathrm{ch}_T(T^{\mathrm{vir}}_{\mathcal{M}_n}) \;=\;
 \bigl(\textstyle\sum_{i=1}^{3} e^{-\eps_i} - 1\bigr)
\]
(three loops supply three tangent contributions with weights
$\eps_1, \eps_2, \eps_3$; three commutator obstructions supply three
factors $(z - w)^{-1}$ in the denominator). The displayed $\omega$ is the unique
$S_2$-symmetric rational function with this $T$-character content. See SV~2013
\S3 for the full derivation.
\end{proof}

\begin{remark}[The $1/(m!n!)$ prefactor]
\label{w17:rem:prefactor}
The prefactor normalises the $S_{m+n}$-sum so that it agrees with the
$(m, n)$-shuffle sum (cosets $S_{m+n}/(S_m \times S_n)$) rather than over-counting
by $m!n!$ (the stabiliser). Either convention determines the same algebraic
structure; the displayed one is that of Schiffmann--Vasserot~2013 and makes the
$m = n = 1$ computation below transparent.
\end{remark}

\subsubsection{Low-order shuffle verification: $p_1 \star p_1$}
\label{w17:subsubsec:p1-star-p1}

\begin{proposition}[Level $m = n = 1$ shuffle product]
\label{w17:prop:p1-star-p1}
\ClaimStatusDerivation\
Let $p_1(z) = z$. Then, using~\eqref{w17:eq:shuffle-product}
at $m = n = 1$:
\begin{equation}
\label{w17:eq:p1-star-p1}
 (p_1 \star p_1)(x_1, x_2) \;=\; x_1 x_2 \cdot
 \bigl[\omega(x_1, x_2) + \omega(x_2, x_1)\bigr]
 \;=\; x_1 x_2 \cdot \left[2 + \frac{2\,\sigma_2}{(x_1 - x_2)^2}\right],
\end{equation}
where $\sigma_2 := \eps_1 \eps_2 + \eps_2 \eps_3 + \eps_3 \eps_1$ is the second
elementary symmetric polynomial. On the CY$_3$ slice
$\eps_1 + \eps_2 + \eps_3 = 0$, the identity
$\sigma_2 = -\tfrac{1}{2}(\eps_1^2 + \eps_2^2 + \eps_3^2)$ holds.
\end{proposition}

\begin{proof}
With $F = p_1(x_1) = x_1$, $G = p_1(x_2) = x_2$, $m = n = 1$, the
prefactor is $1/(1! \cdot 1!) = 1$, and $S_2 = \{e, (1\,2)\}$ so
\[
 (p_1 \star p_1)(x_1, x_2)
 \;=\; \sum_{\sigma \in S_2} \sigma \cdot
 \bigl[x_1 \cdot x_2 \cdot \omega(x_2, x_1)\bigr]
 \;=\; x_1 x_2 \cdot \omega(x_2, x_1) + x_2 x_1 \cdot \omega(x_1, x_2).
\]
Using $x_1 x_2 = x_2 x_1$ gives
$(p_1 \star p_1) = x_1 x_2 \cdot [\omega(x_1, x_2) + \omega(x_2, x_1)]$.

Set $u := x_1 - x_2$. Then from~\eqref{w17:eq:shuffle-kernel}:
\[
 \omega(x_1, x_2) \;=\; \frac{(u - \eps_1)(u - \eps_2)(u - \eps_3)}{u^3},
 \qquad
 \omega(x_2, x_1) \;=\; \frac{(-u - \eps_1)(-u - \eps_2)(-u - \eps_3)}{(-u)^3}.
\]
Pulling out three signs from the numerator of $\omega(x_2, x_1)$ against one
from the denominator gives
$\omega(x_2, x_1) = (u + \eps_1)(u + \eps_2)(u + \eps_3)/u^3$.

Expanding the numerator of the sum:
\begin{align*}
 &(u - \eps_1)(u - \eps_2)(u - \eps_3) + (u + \eps_1)(u + \eps_2)(u + \eps_3) \\
 &\quad = \bigl[u^3 - (\eps_1 + \eps_2 + \eps_3) u^2
 + (\eps_1\eps_2 + \eps_2\eps_3 + \eps_3\eps_1) u - \eps_1\eps_2\eps_3\bigr] \\
 &\quad + \bigl[u^3 + (\eps_1 + \eps_2 + \eps_3) u^2
 + (\eps_1\eps_2 + \eps_2\eps_3 + \eps_3\eps_1) u + \eps_1\eps_2\eps_3\bigr] \\
 &\quad = 2 u^3 + 2 \sigma_2\, u.
\end{align*}
The $u^2$ and constant terms cancel; only the odd-degree terms in $u$ survive.
Dividing by $u^3$ gives $\omega(x_1, x_2) + \omega(x_2, x_1) = 2 + 2\sigma_2/u^2$,
which is~\eqref{w17:eq:p1-star-p1}.

The identity $\sigma_2 = -\tfrac{1}{2}(\eps_1^2 + \eps_2^2 + \eps_3^2)$ on the
CY$_3$ slice follows from
$(\eps_1 + \eps_2 + \eps_3)^2 = \sum \eps_i^2 + 2 \sigma_2 = 0$.
\end{proof}

\begin{verification}[Three paths to Proposition~\ref{w17:prop:p1-star-p1}]
\label{w17:verif:p1-star-p1}\mbox{}
\begin{enumerate}
\item \emph{Direct shuffle expansion.} Above.
\item \emph{Symmetry check.} $p_1 \star p_1$ must lie in
$\mathbb{F}[x_1, x_2]^{S_2}$; the expression $x_1 x_2 [2 + 2\sigma_2/(x_1 - x_2)^2]$
is manifestly $S_2$-symmetric since $(x_1 - x_2)^2$ is.
\item \emph{Scaling check.} Under $\eps_i \mapsto \lambda \eps_i$,
$\sigma_2 \mapsto \lambda^2 \sigma_2$, and the pole term carries weight $\lambda^2$
as expected. Under $x_i \mapsto \lambda x_i$, both the polynomial piece
$x_1 x_2$ and the pole piece $x_1 x_2/(x_1 - x_2)^2$ transform consistently.
The scaling cross-check is a genuine third path.
\end{enumerate}
\end{verification}

\subsubsection{Positive half $Y^+(\gghone)$ and the Drinfeld double}
\label{w17:subsubsec:positive-half}

\begin{definition}[Positive half]
\label{w17:def:positive-half}
\ClaimStatusDefinition\
\[
 Y^+(\gghone) \;:=\; \CoHA(\C^3)_+ \;:=\;
 \bigl\{F \in \Sh : F \text{ has no poles as any } z_i \to z_j\bigr\}.
\]
Equivalently, the sub-shuffle algebra generated by the power sums
$p_k := \sum_i z_i^k$ for $k \geq 0$.
\end{definition}

The pole piece $2 \sigma_2 x_1 x_2/(x_1 - x_2)^2$ in~\eqref{w17:eq:p1-star-p1}
is \emph{not} in $Y^+$ as a function: $Y^+$ is the pole-free closure under the
$\star$-product. In the Drinfeld-currents presentation below, this pole
translates into a commutator relation $[e_0, e_0]$ of the positive Cartan-current
mode.

\begin{theorem}[Tsymbaliuk: Drinfeld-double presentation]
\label{w17:thm:tsymbaliuk-double}
\ClaimStatusProvedElsewhere\ (Tsymbaliuk~2017 \texttt{arXiv:1703.04551} Thm.~1.1.)
The Drinfeld double
\[
 Y(\gghone) \;:=\; Y^+(\gghone) \bowtie Y^0(\gghone) \bowtie Y^-(\gghone)
\]
is isomorphic, as a topological Hopf algebra, to the affine Yangian of $\gghone$
in the Drinfeld-currents presentation: generators
$\{e_k, f_k, \psi_k^\pm\}_{k \geq 0}$ with generating series
\begin{equation}
\label{w17:eq:drinfeld-currents}
 e(u) = \sum_{k \geq 0} e_k u^{-k-1}, \qquad f(u) = \sum_{k \geq 0} f_k u^{-k-1},
 \qquad \psi^\pm(u) = 1 + \sum_{k \geq 0} \psi_k^\pm u^{-k-1},
\end{equation}
and defining relation
\begin{equation}
\label{w17:eq:psi-psi-relation}
 \boxed{\;\psi^\pm(u) \, \psi^\pm(v) \;=\; \varphi(u - v)\,
 \psi^\pm(v) \, \psi^\pm(u)\;}
\end{equation}
with structure function
\begin{equation}
\label{w17:eq:structure-function}
 \varphi(u) \;=\; \frac{(u + \eps_1)(u + \eps_2)(u + \eps_3)}{u^3}.
\end{equation}
The remaining relations (commutation of $e$ with $e$, of $f$ with $f$, of $e$ with
$\psi$, of $f$ with $\psi$, and the Serre-type cubic relations) follow from the
shuffle-algebra structure via the Drinfeld pairing; see Tsymbaliuk~2017 \S3 for
the full list.
\end{theorem}

\begin{remark}[Shuffle kernel vs.\ structure function]
\label{w17:rem:kernel-vs-structure}
The shuffle kernel $\omega(z, w)$ in~\eqref{w17:eq:shuffle-kernel} and the
structure function $\varphi(u)$ in~\eqref{w17:eq:structure-function} are related
by $\omega(z, w) = \varphi(z - w; -\eps_1, -\eps_2, -\eps_3)$: the sign flips
reflect the Drinfeld-pairing-induced inversion when passing from CoHA coordinates
to Cartan-current spectral parameters. The spectral parameter $u$ in $e(u), f(u),
\psi^\pm(u)$ is \emph{not} the same as the moduli-space coordinate $z$ in the
shuffle realisation; conflating the two is AP-CY31.
\end{remark}

\begin{verification}[Three paths to Theorem~\ref{w17:thm:tsymbaliuk-double}]
\label{w17:verif:tsymbaliuk}\mbox{}
\begin{enumerate}
\item \emph{Tsymbaliuk generator-matching.} The proof constructs $Y^+$ from the
SV shuffle, $Y^-$ as the opposite CoHA, and $Y^0$ as the Drinfeld pairing's
Cartan. The coproduct is computed explicitly at low modes and matched with the
affine-Yangian coproduct.
\item \emph{Maulik--Okounkov instanton construction.} Maulik--Okounkov~2012
\texttt{arXiv:1211.1287} construct the Yangian action on equivariant cohomology
of Nakajima quiver varieties; the Jordan quiver (= $\C^3$ after CY$_3$ reduction)
gives exactly $Y(\gghone)$ with structure function~\eqref{w17:eq:structure-function}.
\item \emph{Ding--Iohara--Miki construction.} Ding--Iohara~1997 and Miki~2007
construct the quantum toroidal algebra
$U_{q_1, q_2}(\ddot{\gghone})$; under the limit $q_i = e^{\eps_i}$ and rescaling,
it degenerates to $Y(\gghone)$ with the same structure function. The Miki $S_3$
triality permuting $(\eps_1, \eps_2, \eps_3)$ under CY$_3$ is an independent check
of the symmetry of~\eqref{w17:eq:structure-function}.
\end{enumerate}
\end{verification}

\subsubsection{Two definitions of $\Wilpha$ and their equivalence}
\label{w17:subsubsec:w-infty}

\begin{definition}[$\Wilpha[\lambda]$ \`a la Kac--Radul]
\label{w17:def:winf-kac-radul}
\ClaimStatusDefinition\ (Kac--Radul~1993, \emph{Transform.\ Groups}~\textbf{1}.)
$\Wilpha[\lambda]$ is the unique one-parameter family of vertex operator algebras
(VOAs) characterised by:
\begin{enumerate}
\item central charge $c = \lambda$;
\item exactly one strong generator at each conformal weight $h = 1, 2, 3, \ldots$;
\item OPE structure-constants polynomial in $\lambda$.
\end{enumerate}
At generic $\lambda$, $\Wilpha[\lambda]$ is simple; at $\lambda = N \in \Z_{\geq 0}$
there is a maximal ideal whose quotient is $\mathcal{W}_N \otimes \mathcal{H}$
(an $N$-th $\mathcal{W}$-algebra of principal type tensored with a $\widehat{U(1)}$-
current algebra).
\end{definition}

\begin{definition}[$\Wilpha[N]$ free-field construction]
\label{w17:def:winf-free-field}
\ClaimStatusDefinition\ (Frenkel--Kac--Radul--Wang~1995
\texttt{arXiv:hep-th/9405121}; Awata--Fukuma--Matsuo--Odake~1995.)
Let $N \geq 1$ and let $\mathcal{H}^{\otimes N}$ be the rank-$N$ Heisenberg VOA.
The $\widehat{U(1)}^N$-invariant sub-VOA
\[
 \bigl(\mathcal{H}^{\otimes N}\bigr)^{\widehat{U(1)}^N}
\]
at a specific choice of levels is a strong-generating VOA with one field at each
conformal weight $h = 1, 2, 3, \ldots$. The large-$N$ limit (with appropriate
parameter-rescaling) produces a one-parameter family $\Wilpha[\lambda]$.
\end{definition}

\begin{theorem}[Linshaw: uniqueness of $\Wilpha[\lambda]$]
\label{w17:thm:linshaw}
\ClaimStatusProvedElsewhere\ (Linshaw~2011 \emph{Invariant theory and $\Wilpha$};
monograph \emph{Universal $W$-algebras}.) The Kac--Radul and free-field definitions
determine the same one-parameter family $\Wilpha[\lambda]$:
\begin{enumerate}
\item Both definitions give a simple VOA with central charge $\lambda$ and
generators at weights $1, 2, 3, \ldots$.
\item The OPE structure constants are polynomial in $\lambda$ on both sides.
\item The free-field $\Wilpha[N]$ for $N \in \Z_{\geq 1}$ sits inside the Kac--Radul
family as the truncation at integer $\lambda = N$.
\item Uniqueness of the one-parameter family with stated generating-field content
and polynomial-$\lambda$ OPEs forces equality.
\end{enumerate}
\end{theorem}

\begin{remark}[The $h = 1$ generator is the $\widehat{U(1)}$ current]
\label{w17:rem:h1-generator}
The single field of conformal weight $h = 1$ in Theorem~\ref{w17:thm:linshaw}(1)
is, up to normalisation, the $\widehat{U(1)}$ Heisenberg current $J(z)$; higher-$h$
generators $W_2(z), W_3(z), \ldots$ are normal-ordered composites of $J$ and its
derivatives. The ``$1 +$'' in the name $\Wilpha$ refers to this $h = 1$ generator:
the family is generated by a $\widehat{U(1)}$-current ``plus'' one field at each
$h = 2, 3, \ldots$ ad infinitum.
\end{remark}

\subsubsection{The affine-Yangian $\to \Wilpha$-endomorphism map}
\label{w17:subsubsec:yangian-to-w-end}

\begin{theorem}[Schiffmann--Vasserot + Gaiotto--Rap\v c\'ak evaluation]
\label{w17:thm:yangian-to-w}
\ClaimStatusProvedElsewhere\ (Schiffmann--Vasserot~2013 \texttt{arXiv:1202.2756};
Gaiotto--Rap\v c\'ak~2017 \texttt{arXiv:1703.00982};
Proch\'azka--Rap\v c\'ak~2018 \texttt{arXiv:1807.11304}.) There is a
surjective associative-algebra homomorphism
\begin{equation}
\label{w17:eq:evaluation}
 \mathrm{ev}_\lambda :
 Y(\gghone) \;\twoheadrightarrow\; \mathrm{End}\bigl(\Wilpha[\lambda]\text{-vacuum}\bigr)
\end{equation}
with $\lambda$-dependent kernel. The image realises the affine-Yangian modes as
the mode-algebra acting on the $\Wilpha[\lambda]$-vacuum module; $Y^+$ generates
positive modes, $Y^-$ negative modes, and $Y^0$ the Cartan.
\end{theorem}

\begin{remark}[The three separations, explicit]
\label{w17:rem:three-separations}
Three identifications must be held separate (AP-CY7 and AP934 in
\texttt{notes/antipatterns\_catalogue.md}):
\begin{itemize}
\item $\CoHA(\C^3) \neq Y(\gghone)$: the CoHA is $Y^+$, the positive half only;
the double additionally requires $Y^0$ (Cartan, Drinfeld-pairing) and $Y^-$
(opposite CoHA).
\item $Y(\gghone) \neq \Wilpha[\lambda]$: the Yangian is an associative
$E_1$-algebra; $\Wilpha[\lambda]$ is a vertex algebra (holomorphic factorisation
algebra on $\C$, per Costello--Gwilliam~2017 Vol~I Thm.~5.3.3). The evaluation
$\mathrm{ev}_\lambda$ sends the Yangian to the endomorphism algebra of the
$\Wilpha$-vacuum, not to an isomorphic copy.
\item $\CoHA(\C^3) \neq \Wilpha[\lambda]$ \emph{a fortiori}: the CoHA
is not a vertex algebra. The slogan ``CoHA $= \Wilpha$'' is a character-level
coincidence --- the positive-mode character of $\Wilpha[\lambda]$-vacuum
equals $\chi(\CoHA(\C^3))$ in the $T$-localisation --- not a bar-level
equivalence. AP934 forbids the manifesto-level misidentification.
\end{itemize}
What is unconditional: the triangle holds at the level of mode algebras.
$Y^+(\gghone)$ generates the positive modes of $\Wilpha[\lambda]$ via
$\mathrm{ev}_\lambda$; $Y^-$ generates negative modes; $Y^0$ is the Cartan.
The CoHA is the $Y^+$-half.
\end{remark}

\subsubsection{Two $\lambda$-parameters: $\lambda_{\mathrm{Tr}}$ vs.\ the CY$_3$ slice}
\label{w17:subsubsec:two-lambdas}

\begin{remark}[Truncation parameter vs.\ CY$_3$ identity]
\label{w17:rem:lambda-tr-vs-cy3}
Two parameters coexist and must not be conflated (AP-CY discipline, and
treatise \S``Convention clarification'').
\begin{enumerate}
\item \emph{CY$_3$ slice.} The CY$_3$ condition on torus weights,
$\eps_1 + \eps_2 + \eps_3 = 0$, is an identity on parameters --- a codimension-$1$
slice of the three-parameter $(\eps_1, \eps_2, \eps_3)$-space. On the slice, two
independent equivariant parameters remain (up to overall scale).
\item \emph{Truncation parameter $\lambda_{\mathrm{Tr}}$.} The parameter
controlling $\Wilpha[\lambda]$ on the gauge-theory (Gaiotto--Rap\v c\'ak) side is
\[
 \lambda_{\mathrm{Tr}} \;=\; \frac{\eps_1 + \eps_2}{\eps_3}
 \qquad (\eps_3\text{ a free parameter, not fixed by CY$_3$}).
\]
This is a \emph{ratio} of torus weights.
\end{enumerate}

\emph{CY$_3$ does not collapse the family.} Substituting
$\eps_1 + \eps_2 = -\eps_3$ into $\lambda_{\mathrm{Tr}}$ gives the numerical value
$\lambda_{\mathrm{Tr}} = -1$ at every point of the CY$_3$ slice, but this is an
illegal composition of operations: $\lambda_{\mathrm{Tr}}$ is defined on the
three-parameter domain before the slice; on the slice, the one-parameter family
$\Wilpha[\lambda]$ is parametrised instead by the remaining free ratio
$\eps_1 : \eps_2$, not by $\lambda_{\mathrm{Tr}}$. The canonical slice parameter
is the Gaiotto--Rap\v c\'ak $\lambda_{\mathrm{GR}}$, which depends on
$\eps_1/\eps_2$ and varies continuously on the slice, generating the full
$\Wilpha[\lambda_{\mathrm{GR}}]$-family.

At the further-distinguished point $\eps_1 = \eps_2$ (hence
$\eps_1 = \eps_2 = -\eps_3/2$), $\lambda_{\mathrm{Tr}}$ takes a particular
numerical value, but this is a single point on the slice, not the whole slice.
\end{remark}

\begin{remark}[Gaiotto--Rap\v c\'ak parametrisation]
\label{w17:rem:gr-parametrisation}
The Gaiotto--Rap\v c\'ak parametrisation $\lambda_{\mathrm{GR}}$ encodes the
central charge of $\Wilpha[\lambda_{\mathrm{GR}}]$ as a symmetric triple-rational
functional of $(\eps_1, \eps_2, \eps_3)$:
\[
 c \;=\; 2 + \frac{(\eps_1 + \eps_2 + \eps_3)^2 - (\eps_1^2 + \eps_2^2 + \eps_3^2)
 \cdot \lambda(\lambda - 1)(\lambda - \lambda_{\mathrm{GR}})}{\cdots}
\]
(precise form depends on convention; see Gaiotto--Rap\v c\'ak~2017 \S3). On the
CY$_3$ slice $\eps_1 + \eps_2 + \eps_3 = 0$, the formula simplifies; it does not
reduce to a single numerical value. $\lambda_{\mathrm{GR}}$ and $\lambda_{\mathrm{Tr}}$
are related by a M\"obius transformation in the ambient three-parameter space;
neither is ``the'' $\lambda$.
\end{remark}

\subsubsection{Gauge-theory origin: 4D $\mathcal{N} = 2$ $U(1)$ with $\Omega$-background}
\label{w17:subsubsec:gauge-theory-origin}

\begin{definition}[Setup]
\label{w17:def:gauge-theory-setup}
\ClaimStatusDefinition\ 4D $\mathcal{N} = 2$ pure $U(1)$ gauge theory on
$\R^4 = \C^2_{\eps_1, \eps_2}$ with $\Omega$-background rotating the two complex
planes with weights $\eps_1, \eps_2$. The instanton moduli is
$\bigsqcup_{n \geq 0} \Hilb^n(\C^2)$.
\end{definition}

\begin{theorem}[Nekrasov $U(1)$ partition function]
\label{w17:thm:nekrasov-u1}
\ClaimStatusProvedElsewhere\ (Nekrasov~2003 \texttt{arXiv:hep-th/0206161}.)
\[
 \Zinst^{U(1)}(\eps_1, \eps_2; q)
 \;=\; \sum_{n \geq 0} q^n \cdot \chi_T(\Hilb^n(\C^2))
 \;=\; \prod_{n \geq 1} \frac{1}{1 - q^n},
\]
the 2d-partition generating function (a free-boson Fock-space character).
\end{theorem}

\begin{remark}[Spectral line and $\Wilpha$]
\label{w17:rem:spectral-line}
Adding a holomorphic-topological (HT) spectral line $\C_z$ transverse to
$\C^2_{\eps_1, \eps_2}$, the total spacetime is $\C^2 \times \C_z$. Under the
$\Omega$-twist on $\C^2$ and HT-twist on $\C_z$, the observable algebra is a
factorisation algebra on $\C_z$ with holomorphic locality (Costello--Gwilliam~2017
Vol~I Thm.~5.3.3), equivalent to a vertex algebra.

\ClaimStatusConjectured\ (Costello--Gaiotto~2018 \texttt{arXiv:1810.01970};
Costello--Paquette~2020 \texttt{arXiv:2009.04834}.) This vertex algebra is
$\Wilpha[\lambda]$ with $\lambda$ determined by $(\eps_1, \eps_2, \eps_3)$
via~\eqref{w17:eq:structure-function}. Low-order OPE checks agree; a
single-citation identification theorem is Costello--Paquette~2020.
\end{remark}

\begin{remark}[Nekrasov partition function as $Y^+$-module character]
\label{w17:rem:nekrasov-as-character}
From the Yangian side, $\Zinst^{U(1)}$ is the character of the $Y^+(\gghone)$-module
$\bigoplus_n H^*_T(\Hilb^n(\C^2))$ (Maulik--Okounkov~2012 Thm.~12.1.1;
Schiffmann--Vasserot~2013):
\[
 \chi\bigl(Y^+(\gghone)\text{-mod on } \bigoplus_n H^*_T(\Hilb^n(\C^2))\bigr)
 \;=\; \Zinst^{U(1)}(\eps_1, \eps_2; q).
\]
The three-way pedigree: CoHA $\to$ $Y^+$ (Schiffmann--Vasserot shuffle)
$\to$ Nekrasov partition function (Maulik--Okounkov instanton-moduli character),
with the chiral-algebra face $\mathrm{ev}_\lambda: Y(\gghone) \to
\mathrm{End}(\Wilpha[\lambda]\text{-vac})$ supplying the vertex-algebra incarnation.
\end{remark}

%--------------------------------------------------------------------------------
\subsection{Attack-heal cycles}
\label{w17:subsec:attack-heal-cycles}
%--------------------------------------------------------------------------------

The protocol mandates $\geq 5$ attack-heal cycles targeting the most
conflation-prone steps of the chain; the eight cycles below cover them.

\subsubsection{Attack-heal 1: the Jacobi sign convention}
\label{w17:subsubsec:ah1-jacobi-sign}

\paragraph{Attack.}
The cyclic-derivative computation of Proposition~\ref{w17:prop:jacobi-c3} gives
$\partial_X W = [Y, Z]$, $\partial_Y W = [Z, X]$, $\partial_Z W = [X, Y]$.
A different sign convention (e.g.\ reversing cyclic rotation direction) would
give $[Z, Y]$, $[X, Z]$, $[Y, X]$ instead. Does this change affect
$\Jac(Q, W) = \mathcal{O}_{\C^3}$?

\paragraph{Heal.}
The cyclic derivative $\partial_a: kQ_{\mathrm{cyc}} \to kQ$ is uniquely determined
by the two requirements $\partial_a(b) = \delta_{a, b}$ on generators and
compatibility with cyclic rotation on $kQ_{\mathrm{cyc}}$ (Ginzburg~2006
Def.~3.5.1). This fixes the convention: reversing rotation direction gives a
different operator on $kQ$ that does not descend to $kQ_{\mathrm{cyc}}$ consistently.

Under the fixed cyclic convention, $\partial_X W = [Y, Z]$; reversing would
give $[Z, Y] = -[Y, Z]$. Since the Jacobi ideal is generated by $\partial_a W$,
and $[Z, Y] \in \langle [Y, Z]\rangle$, the ideal is unchanged. The sign
convention controls the \emph{representative} of each generator, not the
generated ideal. Similarly for $\partial_Y, \partial_Z$. Hence
$\Jac(Q, W) = k\langle X, Y, Z\rangle / \langle [X, Y], [Y, Z], [Z, X]\rangle
= k[X, Y, Z] = \mathcal{O}_{\C^3}$, convention-independent.

An equivalent check: the three relations form the ideal of a commutative
algebra, and the commutative quotient is uniquely $k[X, Y, Z]$ regardless of
which signed form of each commutator appears as generator.
\hfill$\square$

\subsubsection{Attack-heal 2: Plane partitions from $T$-fixed commuting triples}
\label{w17:subsubsec:ah2-plane-partitions}

\paragraph{Attack.}
The parametrisation of $\Hilb^n(\C^3)^T$ by plane partitions
(Lemma~\ref{w17:lem:plane-partitions}) uses that a $T$-fixed commuting triple
is necessarily diagonal (monomial). But a generic commuting triple is
\emph{not} diagonal; the commuting variety has non-reduced components with
nilpotent generators. Why does the $T$-fixedness force diagonality?

\paragraph{Heal.}
$T$-fixedness is a strong constraint: the three matrices $X, Y, Z$ must each
lie in a specific $T$-weight space of $\mathrm{End}(\C^n)$, and the common
$T$-eigenspace decomposition must be simultaneous. Concretely, if the
$n$-dimensional representation decomposes under $T$ as
$\C^n = \bigoplus_{(a, b, c) \in S} V_{a, b, c}$ with $V_{a, b, c}$ the
$T$-weight-$(a, b, c)$ subspace, then $X: V_{a, b, c} \to V_{a+1, b, c}$,
$Y: V_{a, b, c} \to V_{a, b+1, c}$, $Z: V_{a, b, c} \to V_{a, b, c+1}$
(the weight of $X$ is $\eps_1$, etc.). For generic torus weights
$(\eps_1, \eps_2, \eps_3)$ linearly independent over $\Q$, each $V_{a, b, c}$ is
at most one-dimensional; the multiset $S$ of weight-vectors labels a basis, and
the matrices become diagonal in this basis.

A non-diagonal commuting triple at a generic point of the commuting variety
(e.g.\ Jordan block) lies away from the $T$-fixed locus precisely because its
spectrum does not split into $T$-weight orbits under the action.

The resulting bijection between $T$-fixed ideals and multisets with
downward-closure is the plane-partition correspondence (MNOP~2006 Prop.~5.1).
\hfill$\square$

\subsubsection{Attack-heal 3: shuffle formula at $m = n = 1$, pole interpretation}
\label{w17:subsubsec:ah3-shuffle-m-n-1}

\paragraph{Attack.}
Proposition~\ref{w17:prop:p1-star-p1} gives
$p_1 \star p_1 = x_1 x_2 [2 + 2\sigma_2/(x_1 - x_2)^2]$. The pole as
$x_1 \to x_2$ is disturbing: does $p_1 \star p_1$ lie in the positive half
$Y^+$ (defined as pole-free) or is it in the localised shuffle $\Sh$?

\paragraph{Heal.}
Two complementary interpretations.

(1) As an element of the localised shuffle algebra $\Sh$, $p_1 \star p_1$
has the displayed form, including the pole at $x_1 = x_2$. The pole is
a genuine feature of the product in $\Sh$; it encodes the $\Omega$-deformation
controlling the non-commutativity of the Yangian.

(2) As a statement about the generator $p_1 \in Y^+$, both copies of $p_1$
are in $Y^+$ (they have no poles), and their product is computed in $\Sh$
rather than in $Y^+$ directly. The relation
\[
 [p_1, p_1] \;=\; p_1 \star p_1 - p_1 \star p_1 \;=\; 0
\]
in $\Sh$ (since $\star$ is commutative at the polynomial level) does not mean
there are no non-trivial commutators in $Y^+$: the non-triviality shows up at
higher shuffles. For $[p_1, p_2]$ with $p_2(z) = z^2$, the pole contribution
becomes a genuine non-vanishing lower-degree polynomial, realising the
Drinfeld-current relations.

A concrete free-field check: at $\sigma_2 = 0$ (e.g.\ $\eps_1 = -\eps_2$,
$\eps_3 = 0$, the two-dimensional degeneration), the pole vanishes and
$p_1 \star p_1 = 2 x_1 x_2$, the free-boson Heisenberg result
$:J(z) J(w): = 2 J(z) J(w) + $ regular for $J(z) = \partial\phi(z)$. The
$\sigma_2 \neq 0$ pole is the $\Omega$-deformation correction.

\emph{Upshot.} The pole in~\eqref{w17:eq:p1-star-p1} is not a pathology; it is
the shuffle-realisation signature of the non-trivial deformation. $Y^+$ as an
abstract algebra (generated by power sums) is pole-free; the product, computed
via the shuffle, is a Laurent expansion around $x_1 = x_2$ whose polar part is
absorbed into the commutator structure.
\hfill$\square$

\subsubsection{Attack-heal 4: uniqueness of the SV shuffle kernel}
\label{w17:subsubsec:ah4-sv-uniqueness}

\paragraph{Attack.}
Theorem~\ref{w17:thm:sv-shuffle} asserts a specific kernel $\omega$. Is this
kernel rigid, or is there a family of equivalent CoHA-kernels? What determines
the particular rational form of $\omega$?

\paragraph{Heal.}
The kernel is determined up to the rescaling $\eps_i \mapsto c \eps_i$
($c \in \mathbb{F}^\times$), which corresponds to a rescaling of the equivariant
parameters. Beyond this, the kernel is pinned uniquely: the virtual-tangent
character of $T^{\mathrm{vir}}\mathcal{M}_n$ is
\[
 \mathrm{ch}_T(T^{\mathrm{vir}}_{[I]}\mathcal{M}_n) \;=\;
 \mathrm{ch}_T(\bigoplus_{i = 1}^{3} V_{\eps_i}) - \mathrm{ch}_T(\bigoplus_{j = 1}^{3}
 V^{\mathrm{obs}}_j) \;=\;
 \sum_{i = 1}^{3} e^{-\eps_i} - 3 + \cdots
\]
where the $-3$ arises from the three commutator obstructions. Exponentiating
to yield the $T$-character of the Euler class of the virtual tangent bundle
and doing a Chern-class expansion gives the rational kernel with numerator
$\prod_{i = 1}^{3}(z - w - \eps_i)$ and denominator $(z - w)^3$.

The SV theorem identifies the CoHA-product with this shuffle product; the
uniqueness of the displayed kernel comes from the rigidity of the virtual-tangent
character on $\mathcal{M}_n$, which is intrinsic to the Jordan triple quiver
with cubic potential and the CY$_3$ torus action. Alternative proofs
(Maulik--Okounkov via Nakajima quiver varieties; Ding--Iohara--Miki via the
quantum toroidal degeneration) arrive at the same kernel up to
$\eps$-rescaling.
\hfill$\square$

\subsubsection{Attack-heal 5: Tsymbaliuk double isomorphism, explicit}
\label{w17:subsubsec:ah5-tsymbaliuk-explicit}

\paragraph{Attack.}
Tsymbaliuk~2017 identifies the Drinfeld double $Y(\gghone) = Y^+ \bowtie Y^0
\bowtie Y^-$ with the affine Yangian of $\gghone$ in Drinfeld-currents
presentation. Is this isomorphism explicit at the level of generators, or
only abstract?

\paragraph{Heal.}
Explicit. On the CoHA / shuffle side, the generators are:
\begin{itemize}
\item $e_k = p_k \in Y^+$: the $k$-th power-sum, witnessing the positive
Cartan-current mode;
\item $f_k \in Y^-$: the opposite CoHA generator;
\item $\psi^\pm_k \in Y^0$: the Cartan generators, obtained from the
Drinfeld pairing
$\langle\cdot, \cdot\rangle: Y^+ \otimes Y^- \to \mathbb{F}$.
\end{itemize}
On the affine-Yangian side, these are the Drinfeld-currents generators of
Theorem~\ref{w17:thm:tsymbaliuk-double}. The $\psi\psi$-commutation
relation~\eqref{w17:eq:psi-psi-relation} is identical on both sides:
\[
 \psi^\pm(u)\, \psi^\pm(v)
 \;=\; \frac{(u - v + \eps_1)(u - v + \eps_2)(u - v + \eps_3)}{(u - v)^3} \cdot
 \psi^\pm(v)\, \psi^\pm(u).
\]

Verifying at low modes: at $k = 0$, $\psi_0^+ = 1$ (normalisation); at
$k = 1$, $\psi_1^+ = \eps_1 + \eps_2 + \eps_3 = 0$ on the CY$_3$ slice
(the leading-$1/u$ coefficient of $\varphi(u - v) - 1$ vanishes). At $k \geq 2$,
the relations become non-trivial and give the explicit Drinfeld-currents
presentation.
\hfill$\square$

\subsubsection{Attack-heal 6: $\Wilpha$ Kac--Radul vs.\ free-field}
\label{w17:subsubsec:ah6-w-inf}

\paragraph{Attack.}
Two definitions of $\Wilpha$ (Kac--Radul's uniqueness-under-strong-generation
characterisation; free-field as $\widehat{U(1)}^N$-invariants) look different.
Why are they equivalent, and why does the equivalence not require a separate
proof at every $N$?

\paragraph{Heal.}
Linshaw's 2011 theorem establishes the equivalence via five steps:
\begin{enumerate}
\item Both definitions yield a VOA with central charge $\lambda$ and strong
generators at weights $h = 1, 2, 3, \ldots$: Kac--Radul by construction,
free-field by the classical invariant theory of the symmetric group $S_N$
acting on symmetric functions of $N$ oscillators.
\item OPE structure constants are polynomial in $\lambda$: Kac--Radul by
definition; free-field by explicit large-$N$ computation in
Linshaw~2011.
\item The free-field $\Wilpha[N]$, $N \in \Z$, sit at integer points
$\lambda = N$ inside the Kac--Radul family.
\item ``Strong reconstruction'' for VOAs (Linshaw monograph \emph{Universal
$W$-algebras}): a VOA with specified strong generators and specified
OPE-closure is determined uniquely up to isomorphism.
\item Matching (1)--(3) via the reconstruction theorem forces equality.
\end{enumerate}
The equivalence does not require a separate proof at every $N$ because the
family is polynomial in $\lambda$ (step (2)); once two analytic families agree
at $\Z$-many points, they are equal. The reconstruction theorem (step (4))
removes any residual freedom.

\hfill$\square$

\subsubsection{Attack-heal 7: $\lambda_{\mathrm{Tr}}$ substitution on the CY$_3$ slice}
\label{w17:subsubsec:ah7-lambda-tr}

\paragraph{Attack.}
$\lambda_{\mathrm{Tr}} = (\eps_1 + \eps_2)/\eps_3$, and CY$_3$ gives
$\eps_1 + \eps_2 + \eps_3 = 0$, so substituting: $\lambda_{\mathrm{Tr}} =
-\eps_3/\eps_3 = -1$. Does this mean $\Wilpha[\lambda_{\mathrm{Tr}} = -1]$
uniformly, collapsing the CY$_3$ family to a single point?

\paragraph{Heal.}
No: the composition ``$\lambda_{\mathrm{Tr}}$ after CY$_3$'' is not a well-formed
evaluation of a one-parameter family of VOAs. Treating it as one is the
specific conflation that AP-CY discipline and the treatise
(\texttt{CoHA\_to\_W\_infty\_treatise.tex} \S``Convention clarification'')
target.

The correct reading: $\lambda_{\mathrm{Tr}}$ is defined on the three-parameter
$(\eps_1, \eps_2, \eps_3)$-space as a \emph{function}. Its value at each point
of the CY$_3$ slice is indeed $-1$. But on the slice, the free parameters are
not $(\eps_1, \eps_2, \eps_3)$; the slice has dimension two (one scalar relation
among three parameters), and any function constant along the slice is simply
an artifact of the parametrisation, not a collapse of the VOA family.

The canonical slice parameter (Gaiotto--Rap\v c\'ak's $\lambda_{\mathrm{GR}}$)
depends on the ratio $\eps_1/\eps_2$ and varies continuously on the slice,
correctly parametrising the full $\Wilpha[\lambda_{\mathrm{GR}}]$-family.
$\lambda_{\mathrm{Tr}}$ and $\lambda_{\mathrm{GR}}$ are related by a M\"obius
transformation in the ambient three-parameter space; neither is ``the'' $\lambda$.

\emph{Upshot.} The CY$_3$ identity is a constraint on parameters, not a
distinguished numerical evaluation. Treating it as the latter is the AP-CY
conflation. The correct reparametrisation on the slice preserves the
one-parameter family.
\hfill$\square$

\subsubsection{Attack-heal 8: Nekrasov $U(1)$ as $Y^+$-module character}
\label{w17:subsubsec:ah8-nekrasov}

\paragraph{Attack.}
The Nekrasov $U(1)$ partition function is
$\prod_n (1 - q^n)^{-1}$ (2d partition generating function), not the MacMahon
function $\prod_n (1 - q^n)^{-n}$ (3d partition generating function).
Is this the character of a $Y^+(\gghone)$-module?

\paragraph{Heal.}
Yes. The $U(1)$-instanton moduli on $\C^2 = \C^2_{\eps_1, \eps_2}$ is
$\bigsqcup_n \Hilb^n(\C^2)$, and $T^2$-localisation via
Atiyah--Bott gives
\[
 Z^{U(1)}_\Omega(\eps_1, \eps_2; q)
 \;=\; \sum_n q^n \cdot \sum_{[I] \in \Hilb^n(\C^2)^{T^2}}
 \frac{1}{e^{T^2}(T_{[I]}\Hilb^n)}
 \;=\; \prod_n (1 - q^n)^{-1}
\]
in the Abelian sector, by a direct computation at each fixed point (the
equivariant-Euler contribution cancels against the $q$-weight of the partition).

Maulik--Okounkov~2012 Thm.~12.1.1 identifies $Y(\gghone)$ with the instanton
algebra acting on $\bigoplus_n H^*_{T^2}(\Hilb^n(\C^2))$, making this space a
vacuum module for $Y^+(\gghone)$. Its character in $q$ is precisely
$Z^{U(1)}_\Omega$.

The 3d connection: $\Hilb^n(\C^2)$ is the ``flat slice'' of $\Hilb^n(\C^3)$ cut
out by fixing the third coordinate. The 3d MacMahon function $M(q)$ counts
$T^3$-fixed points on $\Hilb^n(\C^3)$; the 2d specialisation
$\prod_n (1 - q^n)^{-1}$ counts $T^2$-fixed points on $\Hilb^n(\C^2)$, a
sublocus. Both are characters of $Y^+$-modules (restricted along the
dimensional specialisation), and the $\C^3$-CoHA Schiffmann--Vasserot shuffle
acts on both via the displayed kernel.
\hfill$\square$

%--------------------------------------------------------------------------------
\subsection{Example 2: resolved conifold}
\label{w17:subsec:resolved-conifold}
%--------------------------------------------------------------------------------

\subsubsection{Geometry and quiver}
\label{w17:subsubsec:conifold-geometry}

\begin{definition}[Resolved conifold]
\label{w17:def:resolved-conifold}
\ClaimStatusDefinition\
$\bY := \mathrm{Tot}\bigl(\mathcal{O}_{\mathbb{P}^1}(-1) \oplus
\mathcal{O}_{\mathbb{P}^1}(-1) \to \mathbb{P}^1\bigr)$,
the total space of two degree-$-1$ line bundles over $\mathbb{P}^1$. By
adjunction $K_{\bY} = K_{\mathbb{P}^1} \otimes \det(\mathcal{O}(-1) \oplus
\mathcal{O}(-1))^{-1} = \mathcal{O}(-2) \otimes \mathcal{O}(2) = \mathcal{O}$,
so $\bY$ is CY$_3$. It is a local (non-compact) toric threefold; its toric fan
has two 3d cones meeting along a compact 2d wall dual to $\mathbb{P}^1$.
\end{definition}

\begin{definition}[Klebanov--Witten quiver $\Qcon$]
\label{w17:def:kw-quiver}
\ClaimStatusDefinition\ Two vertices $\circ, \bullet$ (corresponding to the two
toric fixed points on $\mathbb{P}^1$). Four arrows:
$a_1, a_2: \circ \to \bullet$ and $b_1, b_2: \bullet \to \circ$. The path
algebra $k\Qcon$ is generated by paths in this quiver.
\end{definition}

\begin{definition}[Klebanov--Witten potential]
\label{w17:def:kw-potential}
\ClaimStatusDefinition\
$W_{\mathrm{KW}} := \mathrm{tr}(a_1 b_1 a_2 b_2 - a_1 b_2 a_2 b_1)
\in k\Qcon_{\mathrm{cyc}}$.
\end{definition}

\begin{theorem}[Klebanov--Witten / Szendr\H{o}i: Jacobi $\cong$ conifold NCCR]
\label{w17:thm:conifold-jacobi}
\ClaimStatusProvedElsewhere\ (Klebanov--Witten~1998
\texttt{arXiv:hep-th/9807080}; Szendr\H{o}i~2008 \emph{Geom.\ Topol.}~\textbf{12}
\texttt{arXiv:0705.3419}.) The Jacobi algebra
$\Jac(\Qcon, W_{\mathrm{KW}})$ is isomorphic as a $k$-algebra to the
non-commutative crepant resolution $\Lambda_{\bY}$ of the conifold singularity,
and $D^b(\mathrm{mod}\,\Lambda_{\bY}) \simeq D^b(\mathrm{Coh}\,\bY)$.
\end{theorem}

\begin{remark}[Cyclic derivatives of the KW potential]
\label{w17:rem:kw-cyclic-derivatives}
Direct computation:
\begin{align*}
 \partial_{a_1} W_{\mathrm{KW}} &= b_1 a_2 b_2 - b_2 a_2 b_1, \\
 \partial_{a_2} W_{\mathrm{KW}} &= b_2 a_1 b_1 - b_1 a_1 b_2, \\
 \partial_{b_1} W_{\mathrm{KW}} &= a_2 b_2 a_1 - a_1 b_2 a_2, \\
 \partial_{b_2} W_{\mathrm{KW}} &= a_1 b_1 a_2 - a_2 b_1 a_1,
\end{align*}
imposing ``super-commutativity'' of the quiver generators up to an explicit
pattern that matches the conifold NCCR (Van den Bergh~2004
\texttt{arXiv:math/0207170}).
\end{remark}

\subsubsection{CoHA of the conifold and toroidal shuffle}
\label{w17:subsubsec:conifold-coha}

\begin{definition}[CoHA of the conifold]
\label{w17:def:conifold-coha}
\ClaimStatusDefinition\ For dimension vector $(d_1, d_2) \in \Z_{\geq 0}^2$, the
representation moduli is
$\mathcal{M}_{d_1, d_2} = [\mathrm{Rep}(\Qcon, d_1, d_2)/\prod GL_{d_i}]$, and
\[
 \CoHA(\bY) \;:=\; \bigoplus_{(d_1, d_2)} H^*_T\bigl(\mathcal{M}_{d_1, d_2},
 \phi_{W_{\mathrm{KW}}}\bigr).
\]
\end{definition}

\begin{theorem}[Davison / Morrison--Mozgovoy--Nagao--Szendr\H{o}i: conifold CoHA]
\label{w17:thm:conifold-coha}
\ClaimStatusProvedElsewhere\ (Davison~2012 \texttt{arXiv:1209.4620} for the
explicit identification; Morrison--Mozgovoy--Nagao--Szendr\H{o}i~2010 for
DT-invariants.)
\[
 \CoHA(\bY) \;\cong\; U(\gBPS(\bY))
\]
as associative algebras, where $\gBPS(\bY)$ is a specific Lie algebra determined
by the conifold geometry (primitive-cohomology subspaces of the DT moduli).
\end{theorem}

\begin{theorem}[Rap\v c\'ak--Soibelman--Yang--Zhao: conifold Yangian]
\label{w17:thm:rsyz-conifold}
\ClaimStatusProvedElsewhere\ (Rap\v c\'ak--Soibelman--Yang--Zhao~2020
\texttt{arXiv:2001.10549} Thm~B; Kapranov--Vasserot~2018
\texttt{arXiv:1802.07988} for the K-theoretic toroidal incarnation.)
The positive-half CoHA of the conifold matches the positive half of the
quantum toroidal algebra $U_{q_1, q_2}(\ddot{\mathfrak{gl}}_2)$ at a specific
point of its parameter space, where $q_i = e^{\eps_i}$ and the two equivariant
parameters come from the two non-compact directions.
\end{theorem}

\begin{definition}[Shuffle kernel for conifold]
\label{w17:def:conifold-shuffle-kernel}
\ClaimStatusDefinition\ (Negu\c t~2015 \texttt{arXiv:1505.01528}.)
Over $\mathbb{F}[\eps_1, \eps_2]$, the conifold shuffle kernel is
\[
 \omega_{\mathrm{con}}(z, w) \;:=\;
 \frac{(z - w + \eps_1)(z - w + \eps_2)}{(z - w)(z - w + \eps_1 + \eps_2)}.
\]
The CY$_3$ condition is $\eps_1 + \eps_2 = 0$ on the fibre directions (the
$\mathbb{P}^1$-direction is compact and does not enter the equivariant
parameter), so under CY$_3$, $\omega_{\mathrm{con}}$ simplifies but does not
degenerate.
\end{definition}

\begin{remark}[Why two shuffle kernels differ in structure]
\label{w17:rem:kernel-difference}
The $\C^3$ kernel $\omega(z, w) = \prod_{i = 1}^{3}(z - w - \eps_i)/(z - w)^3$
has three equivariant parameters (three loops at one vertex); the conifold
kernel has two (two pairs of arrows between two vertices). The structural
difference reflects the different quiver shape, not a different formalism:
both are virtual-tangent characters of the respective representation-moduli
stacks. For general toric CY$_3$ without compact 4-cycles, the Schiffmann
kernel is a rational function of $(z - w)$ with poles and zeros controlled
by the quiver dimension vectors.
\end{remark}

\subsubsection{Conifold chiral Yangian and gauge-theory origin}
\label{w17:subsubsec:conifold-chiral}

\begin{remark}[6D hCS on the conifold and defect factorisation algebras]
\label{w17:rem:hcs-on-conifold}
\ClaimStatusConjectured\ (Costello--Li~2016 \texttt{arXiv:1606.00365};
Costello--Gwilliam Vol II \S6 conifold example.) Holomorphic Chern--Simons
(hCS) on $\bY$ supports a factorisation algebra
$\mathcal{F}_{\hCS}(\bY)$; restricting to a defect line $\C_z \subset \bY$
along a generic fibre direction, the defect factorisation algebra
$\mathcal{F}_{\mathrm{defect, con}}$ is a vertex algebra.

\ClaimStatusConjectured\ $\mathcal{F}_{\mathrm{defect, con}}$ is a quantum toroidal
algebra at specific parameters --- equivalently, a deformation of
$\Wilpha[\lambda]$ corresponding to the conifold topology
(Awata--Feigin--Shiraishi~2011 \texttt{arXiv:1112.6074}).

\ClaimStatusOpen\ \emph{Explicit vertex-algebra generators and OPEs of the
conifold chiral Yangian.} Costello--Li~2016 and Costello--Paquette~2020 provide
the framework; explicit formulas for all generators and OPEs do not reduce to
a single primary citation at present.
\end{remark}

\begin{remark}[Two-stage factorisation $\Phi_3$ on the conifold]
\label{w17:rem:phi3-conifold}
Via the Vol~III two-stage functor (per CLAUDE.md):
$\Phi_3 = \SpCh_{\Sigma_2, C} \circ \PhiFA_3$ applied to
$D^b(\mathrm{Coh}\,\bY)$ yields, at Stage~1, the $E_3$-holomorphic factorisation
algebra of polyvector fields on $\bY$ ($\Omega$-deformed via the toric
equivariance); at Stage~2, specialisation to a chosen curve $C \subset \bY$
via integration over $\Sigma_2$ gives the $E_1$-chiral shadow on $C$. The
defect factorisation algebra $\mathcal{F}_{\mathrm{defect, con}}$ on a generic
fibre line $\C_z$ is the case $\Sigma_2 = \mathbb{P}^1_{\mathrm{compact}}$
and $C = \C_z$. AP929 forbids framing $\Phi_3$ as a single-stage functor;
the two-stage realisation is essential.
\end{remark}

%--------------------------------------------------------------------------------
\subsection{Example 3: $K3 \times E$}
\label{w17:subsec:k3-times-e}
%--------------------------------------------------------------------------------

\subsubsection{Setup and obstructions to the $\C^3$ strategy}
\label{w17:subsubsec:k3e-setup}

$X := K3 \times E$: compact smooth CY$_3$ of complex dimension $3$ with
Hodge diamond
\[
 h^{0, 0} = 1, \quad h^{1, 0} = h^{0, 1} = 1, \quad h^{1, 1} = 21 + 1 = 22,
 \quad h^{2, 0} = h^{0, 2} = 1,
\]
where the split reflects the K\"unneth decomposition of $K3$ Hodge
structure~$(1, 20, 1)$ in degree $2$ with $E$ Hodge structure~$(1, 1)$ in
degree~$1$.

\emph{Obstructions to the $\C^3$ localisation strategy:}
\begin{enumerate}
\item $X$ is \emph{not toric}. Generic $K3$ has no non-trivial $\C^\times$-action
(the symplectic holomorphic 2-form $\omega_K$ would be killed by any
$\C^\times$-contraction). Hence equivariant cohomology with respect to a
would-be $T^3$ is ordinary cohomology, and localisation reduces to trivial
identities.
\item $X$ does not admit a quiver-with-potential presentation in the toric
sense; the CoHA construction via $T^3$-equivariant vanishing cycles does not
apply directly.
\item The $E$-factor contributes an extra $H^*(E) = \C[0] \oplus \C[1]$ Hodge
component, introducing a ``rigidity'' into the CY$_3$ structure absent on
local toric CY$_3$.
\end{enumerate}

\subsubsection{Positive effective geometry via reduced DT}
\label{w17:subsubsec:k3e-reduced-dt}

\begin{theorem}[Positive effective geometry on $K3 \times E$]
\label{w17:thm:k3e-reduced-dt}
\ClaimStatusProvedElsewhere\ (Pandharipande--Thomas~2014;
Oberdieck--Pixton~2017 \texttt{arXiv:1706.10100}, \emph{Invent.\ Math.}~\textbf{222}
(2020).) Let $\gamma \in H^{\mathrm{ev}}(X, \Z)$ restrict to a primitive curve
class on the $K3$-factor. The positive-effective moduli is the
Pandharipande--Oberdieck reduced DT stack
\[
 \Meff(X) \;=\; \Mred(X; \gamma\text{ primitive on }K3),
\]
equivariant for the rigid two-factor group
\[
 G_{K3 \times E} \;=\; \C^\times_E \times \Auts(K3),
\]
with $\C^\times_E$ the translation $\C^\times$ on the elliptic factor and
$\Auts(K3) \subset M_{23}$ the symplectic-automorphism group of a Mukai-lattice
polarisation. The reduced DT partition function is
\[
 Z^{\mathrm{red}}_{DT}(X; q, t, p) \;=\;
 -\frac{1}{\Phi_{10}(q, t, p)} \;=\; -\Delta_5^{-2}(q, t, p),
\]
where $\Phi_{10}$ is the Igusa cusp form and $\Delta_5$ is the Gritsenko lift of
$\phi_{0, 1}$ (Gritsenko--Nikulin~1998 \emph{Amer.\ J.\ Math.}~\textbf{119}).
The associated Borcherds--Kac--Moody weight is
\[
 \kappa_{\mathrm{BKM}}(\Delta_5) \;=\; \frac{c_N(0)}{2}\bigg|_{N = 1}
 \;=\; \frac{10}{2} \;=\; 5.
\]
\end{theorem}

\begin{remark}[Why toric localisation is replaced by reduced DT]
\label{w17:rem:toric-replaced-by-reduced}
The $\C^3$ toric localisation strategy (Maulik--Okounkov, Schiffmann--Vasserot)
uses $T$-fixed-point reduction to plane partitions. On $K3 \times E$ the
replacement comprises three inputs (per Maulik--Pandharipande--Thomas~2010
\texttt{arXiv:1001.2719} and Oberdieck--Pixton):
\begin{enumerate}
\item Reduced virtual class: $E$-symplectic form trivialises one copy of the
obstruction bundle, giving a virtual cycle of dimension one higher than the
naive DT class.
\item Residual $\C^\times_E$ via translation on the elliptic factor.
\item Discrete $\Auts(K3) \subset M_{23}$ orbit via symplectic automorphisms
(Mukai~1988). The combined group $G_{K3 \times E}$ acts on moduli, replacing
$T^3$.
\end{enumerate}
Together they supply the equivariance to mimic the toric-localisation strategy
for CoHA construction.
\end{remark}

\subsubsection{Davison BPS Lie algebra: construction and graded-dimension theorem}
\label{w17:subsubsec:k3e-davison}

\begin{definition}[Davison BPS Lie algebra, conditional construction]
\label{w17:def:davison-bps}
\ClaimStatusConjectured\ (Davison~2017 \texttt{arXiv:1601.02479};
Davison--Meinhardt~2020 \emph{Invent.\ Math.}~\textbf{221}.)
\emph{Scope hypothesis:} $X$ a smooth proper CY$_3$ admitting a global
critical-locus / potential presentation for its sheaf moduli (not automatic
for every compact CY$_3$; Davison's construction is formulated for
quiver-with-potential models and their direct geometric counterparts, with the
global critical-chart hypothesis load-bearing).
Under this hypothesis,
\[
 \CoHA^{\mathrm{BPS}}(X) \;:=\;
 H^*\bigl(\mathcal{M}(X), \phi_W \cdot \mathrm{IC}_{\mathcal{M}(X)}\bigr)
\]
is an associative algebra, and there exists a BPS Lie sub-algebra
$\gBPS(X) \subset \CoHA^{\mathrm{BPS}}(X)$ of primitive elements with PBW
theorem $\CoHA^{\mathrm{BPS}}(X) = U(\gBPS(X))$.
\end{definition}

\begin{theorem}[Graded-dimension identification on $K3 \times E$]
\label{w17:thm:k3e-graded-dim}
\ClaimStatusProvedElsewhere\
(Oberdieck--Pixton~2017; Davison--Meinhardt~2020; Borcherds~1998
\emph{Invent.\ Math.}~\textbf{132}; Gritsenko--Nikulin~1998.)
For every $\gamma$ primitive on $K3$,
\[
 \dim \gBPS_{, \gamma}(K3 \times E)
 \;=\; \dim \gDelta_{, \alpha_\gamma}
 \;=\; \mathrm{mult}_{\mathrm{BKM}}(\alpha_\gamma),
\]
where $\gDelta$ is the Borcherds--Kac--Moody superalgebra of Gritsenko--Nikulin
attached to the Borcherds product $\Delta_5$, and $\alpha_\gamma$ is the root of
$\gDelta$ corresponding to $\gamma$ via the Mukai-lattice identification.
The identity unconditionally couples three independent inputs:
\begin{itemize}
\item Oberdieck--Pixton reduced-DT formula: $Z^{\mathrm{red}}_{DT} = -1/\Phi_{10}$;
\item Davison--Meinhardt: $\Omega(\gamma) = \dim \gBPS_{, \gamma}$ where $\Omega$
is the numerical BPS invariant;
\item Borcherds denominator formula: $\Phi_{10} = \Delta_5^2$ expanded as a
product over positive roots of $\gDelta$ with multiplicities.
\end{itemize}
\end{theorem}

\subsubsection{The bracket-level identification: open problem}
\label{w17:subsubsec:k3e-open}

\begin{remark}[Status of the bracket-level identification]
\label{w17:rem:k3e-open-status}
\ClaimStatusOpen\ \emph{Lie-algebra isomorphism}
$\gBPS(K3 \times E) \stackrel{?}{\cong} \gDelta$. Two Lie algebras with equal
weight-space dimensions are not automatically isomorphic: one must check that
the brackets $[\cdot, \cdot]$ coincide, i.e.\ for $\gamma_1 + \gamma_2 = \gamma$:
\[
 [\gBPS_{, \gamma_1}, \gBPS_{, \gamma_2}] \;\to\; \gBPS_{, \gamma}
 \quad\text{vs.\ }\quad
 [\gDelta_{, \alpha_{\gamma_1}}, \gDelta_{, \alpha_{\gamma_2}}]
 \;\to\; \gDelta_{, \alpha_\gamma}.
\]
The first is determined by the CoHA correspondence product; the second by the
BKM denominator formula's root-system structure. Their equality is the open
content of the identification.

\ClaimStatusConjectured\ The identification is \emph{consistent} with
Kontsevich--Soibelman motivic wall-crossing (bracket of BPS states at a
wall = bracket of BKM roots), but a reduction to a single-citation theorem
is not in hand.

The graded-dimension matching
(Theorem~\ref{w17:thm:k3e-graded-dim}) is unconditional; the bracket-level
identification is tracked as AP-CY34 in the catalogue.
\end{remark}

\subsubsection{Chiral BKM on $K3 \times E$: what would it look like?}
\label{w17:subsubsec:k3e-chiral}

\begin{remark}[Conjectural class-$\mathcal{S}$ / M5-brane origin]
\label{w17:rem:k3e-chiral-conjecture}
\ClaimStatusConjectured\ By analogy with the $\C^3$ case: a 4D $\mathcal{N} = 2$
theory obtained from M5-branes on $K3 \times E$ (the Vafa--Witten / $(2, 0)$
setup), with $\Omega$-deformation on $\C^2_{\eps_1, \eps_2}$ transverse to
$K3 \times E$, compactified along a defect curve $\Sigma \subset K3 \times E$
(candidates: the $E$-factor, a hyperelliptic curve in $K3$, or the holomorphic
spectral line after HT-twist), would supply a factorisation algebra on $\Sigma$.

Under the Vol~III $\Phi_3 = \SpCh_{\Sigma_2, C} \circ \PhiFA_3$ applied at
$(K3 \times E, \Sigma)$, the $E_1$-chiral shadow on the reference curve $\Sigma$
is conjectured to be a chiral BKM whose primitive / BPS sub-object is $\gDelta$.

\ClaimStatusOpen\ \emph{Identification of the chiral BKM.} Neither the explicit
chiral-algebra generators, nor the OPE structure, nor the identification with
$\gDelta$ as Lie algebras, are in hand. This is the central conjecture of the
Vol~III $K3 \times E$ programme.
\end{remark}

\subsubsection{Cross-consistency with the treatise}
\label{w17:subsubsec:k3e-treatise-consistency}

The treatise \texttt{CoHA\_to\_W\_infty\_treatise.tex}
(\texttt{\textasciitilde /calabi-yau-quantum-groups/notes/}) enumerates the
same obstructions and status markers for $K3 \times E$; the only technical
refinement is the \texttt{wn:rem:K3xE-open-problem-reconciliation} clause,
which separates the bracket-level and graded-dimension claims explicitly. This
note adopts the same separation:
\begin{itemize}
\item Graded dimensions (Theorem~\ref{w17:thm:k3e-graded-dim}): unconditional.
\item Bracket-level identification
(Remark~\ref{w17:rem:k3e-open-status}): open, consistent with KS wall-crossing
but not reduced to a single citation.
\item Chiral-algebra identification
(Remark~\ref{w17:rem:k3e-chiral-conjecture}): conjectural, absent full
$\Phi_3$-construction at chain level on $K3 \times E$.
\end{itemize}

No inconsistency with the treatise is identified at the level of theorem
statements, scope tags, or open-problem boundaries. The cross-consistency
pass is clean.

%--------------------------------------------------------------------------------
\subsection{Summary table and overall status}
\label{w17:subsec:summary-table}
%--------------------------------------------------------------------------------

\begin{center}
\begin{tabular}{l|l|l|l}
\textbf{Step} & \textbf{$\C^3$} & \textbf{Conifold} & \textbf{$K3 \times E$} \\\hline
Quiver$+$potential & Jordan triple, $\mathrm{tr}(X[Y, Z])$
& 2-vertex / 4-arrow, $W_{\mathrm{KW}}$ & N/A (compact, non-toric) \\
$\Jac \cong \mathcal{O}_{\mathrm{CY}_3}$ & Proposition~\ref{w17:prop:jacobi-c3}
& Klebanov--Witten / Szendr\H{o}i
& N/A \\
CoHA as associative & Kontsevich--Soibelman~2008
& Davison~2012 & Davison~2017 (scope hyp.) \\
Shuffle kernel & Theorem~\ref{w17:thm:sv-shuffle}
& Negu\c t~2015 & open \\
Positive half $Y^+$ & $Y^+(\gghone)$ shuffle
& conifold toroidal $Y^+$ & conjectural ``$Y^+_{K3 \times E}$'' \\
Drinfeld double & Tsymbaliuk~2017 & RSYZ~2020 / KV~2018 & open \\
Identified with VOA & $\Wilpha[\lambda]$-vac (via $\mathrm{ev}_\lambda$)
& quantum toroidal (Costello--Li~2016) & open \\
Gauge-theory origin & 4D $\mathcal{N} = 2$ $U(1)$ $\Omega$-bg
& idem (non-compact) & class-$\mathcal{S}$ $(A_1, \Sigma)$ \\
hCS factorisation alg & Costello--Gwilliam~2017 & Costello--Li~2016 & open \\
BPS Lie algebra & $\gghone$ & toroidal $\ddot{\mathfrak{gl}}_2$
& $\gDelta$ (graded-dim only) \\
\end{tabular}
\end{center}

\emph{Interpretation.} The first column is rigid and fully derived: every step
between $\Jac \cong \mathcal{O}_{\C^3}$ and $Y(\gghone) \to \mathrm{End}
(\Wilpha[\lambda]\text{-vac})$ is proved, cited, and verified at low orders.
The second column is partially rigid: the CoHA construction works on the
conifold; the Drinfeld-double / toroidal identification is partial
(Rap\v c\'ak--Soibelman--Yang--Zhao / Kapranov--Vasserot); the chiral-algebra
face via Costello--Li is conjectural at the explicit-generator level. The
third column is the programme's target: the graded-dimension matching is
unconditional, the bracket-level identification is open, and the
chiral-algebra incarnation is conjectural.

The triangle $\CoHA(\C^3) = Y^+(\gghone) \to \mathrm{End}
(\Wilpha[\lambda]\text{-vac})$ is the template that the $K3 \times E$
programme targets; AP934 reminds us that ``$\CoHA(\C^3) = \Wilpha$'' is the
character-level coincidence, not the bar-level equivalence, and
$\CoHA(\C^3) = Y^+(\gghone)$ is the full quantum identification.

%--------------------------------------------------------------------------------
\subsection{The three load-bearing separations revisited}
\label{w17:subsec:three-separations-final}
%--------------------------------------------------------------------------------

Per Remark~\ref{w17:rem:three-separations} and AP-CY7 / AP934, each of the
following identifications must be held separate at every reader-facing citation:
\begin{enumerate}
\item $\boxed{\CoHA(\C^3) \neq Y(\gghone).}$ CoHA is the positive half $Y^+$;
the Drinfeld double further requires $Y^0$ (Cartan, from Drinfeld pairing) and
$Y^-$ (opposite CoHA). The identity $\CoHA(\C^3) = Y^+(\gghone)$ holds;
$\CoHA(\C^3) = Y(\gghone)$ does not.
\item $\boxed{Y(\gghone) \neq \Wilpha[\lambda].}$ The Yangian is an associative
$E_1$-algebra; $\Wilpha[\lambda]$ is a vertex algebra (holomorphic factorisation
algebra on $\C$). The evaluation $\mathrm{ev}_\lambda: Y(\gghone) \to
\mathrm{End}(\Wilpha[\lambda]\text{-vac})$ realises the Yangian as the mode
algebra on the vacuum module; it is \emph{not} an isomorphism.
\item $\boxed{\CoHA(\C^3) \neq \Wilpha[\lambda].}$ The CoHA is associative;
$\Wilpha$ is a vertex algebra. Their equality is a character-level
coincidence (graded dimensions match under $\mathrm{ev}_\lambda$-image), not a
bar-level equivalence (no isomorphism of underlying algebraic structures).
AP934 is the manifesto-level version of this separation: ``CoHA$(\C^3) =
\Wilpha$'' is the slogan; the correct identification is $\CoHA(\C^3) =
Y^+(\gghone)$.
\end{enumerate}

What is unconditional: the three objects sit on the triangle
$\CoHA(\C^3) = Y^+(\gghone) \to \mathrm{End}(\Wilpha[\lambda]\text{-vac})$,
with $Y^+$ generating positive modes, $Y^-$ generating negative modes, $Y^0$
the Cartan, and $\mathrm{ev}_\lambda$ the evaluation on the vacuum module.

%--------------------------------------------------------------------------------
\subsection{Cross-volume anchor}
\label{w17:subsec:cross-volume-anchor}
%--------------------------------------------------------------------------------

Per CLAUDE.md's two-stage factorisation framework (AP929), this note's
reference-curve Yangian statements on $\C^3$, on the conifold, and
(conjecturally) on $K3 \times E$, are Stage-2 $E_1$-chiral shadows of Vol~III's
$\PhiFA_3$-assignment on the respective CY$_3$. Vol~III's two-stage
factorisation $\Phi_3 = \SpCh_{\Sigma_2, C} \circ \PhiFA_3$ is the target
object; this note's chain-level identifications are the Stage-2 ($E_1$-chiral)
shadow of Vol~III's Stage-1 ($E_3$-holomorphic factorisation algebra) output
for each CY$_3$.

Specifically:
\begin{itemize}
\item On $\C^3$: Stage 1 gives the polyvector-field $E_3$-algebra on $\C^3$
with $\Omega$-deformation parameter $(\eps_1, \eps_2, \eps_3)$ (toric
equivariance replaces Costello--Gwilliam--Li locality in the local toric
regime). Stage 2 specialises to $C = \C_z$ (the spectral line) via integration
over $\Sigma_2 = \C^2_{\eps_1, \eps_2}$, landing in the $E_1$-chiral shadow
$\Wilpha[\lambda]$. The $Y^+(\gghone) \subset \CoHA(\C^3)$ is the mode
algebra's positive half.
\item On $\bY$ (conifold): Stage 1 gives polyvector fields on $\bY$ with
$\Omega$-deformation from the conifold torus $T^2$. Stage 2 specialises to
a generic fibre curve; the $E_1$-chiral shadow is the (conjectural) conifold
chiral Yangian.
\item On $X = K3 \times E$: Stage 1 is conjectural at chain level
(global-critical-chart hypothesis of Davison); Stage 2's $E_1$-chiral shadow is
the conjectural chiral BKM on the reference curve $C \subset X$.
\end{itemize}

AP929 forbids framing $\Phi_3$ as a single-stage functor; the two-stage
decomposition preserves the $(\Sigma_2, C)$-family structure of the
$E_1$-chiral shadows. A single CY$_3$ category admits a family of $E_1$-chiral
shadows indexed by $(\Sigma_2, C)$; Theorem~A governs each shadow on its curve.

%--------------------------------------------------------------------------------
\subsection{Closing: the programme-scoping rule}
\label{w17:subsec:programme-scoping}
%--------------------------------------------------------------------------------

The $\C^3$-triangle
$\CoHA(\C^3) = Y^+(\gghone) \to \mathrm{End}(\Wilpha[\lambda]\text{-vac})$
is the rigid foundation for the first column of the Vol~III programme matrix.
The programme targets the third column ($K3 \times E$) and its conjectural
triangle $\CoHA(K3 \times E) \cong$ ``$Y^+_{K3 \times E}$'' $\cong$
chiral-$\Delta_5$-BKM, which at present is at the level of graded-dimension
matching only (Theorem~\ref{w17:thm:k3e-graded-dim}).

The $\C^3$-triangle is the template: its unconditional content ---
Proposition~\ref{w17:prop:jacobi-c3}, the SV shuffle
(Theorem~\ref{w17:thm:sv-shuffle}), the Tsymbaliuk double
(Theorem~\ref{w17:thm:tsymbaliuk-double}), the Linshaw $\Wilpha$ uniqueness
(Theorem~\ref{w17:thm:linshaw}), the $\mathrm{ev}_\lambda$ surjection
(Theorem~\ref{w17:thm:yangian-to-w}) --- is the target of the
$K3 \times E$ programme, obstructed by the non-toric nature of $K3$ and the
absence of continuous $\C^\times$-action. The conifold is the intermediate
case: partial rigid construction (Davison~2012; RSYZ~2020) with conjectural
vertex-algebra face (Costello--Li~2016).

\emph{The programme's centre of gravity} is the last row of the last column:
the conjectural identification
$\gBPS(K3 \times E) \cong \gDelta$ as Borcherds--Kac--Moody Lie superalgebras,
and the associated chiral BKM vertex algebra realised via the conjectural
$\Phi_3$ on the $K3 \times E$ category.

\emph{End of note.}
